% The integrated commentary of the concise study, beginning on physical page
% 5. It is built on questions that more than one of the three interpretations
% answers: why the servants may not weed and how far correction may go; what
% the servants do meanwhile; why a God who has promised must be asked; what
% the Gradual's thanks and the Collect's trust say of those who sing and pray
% them; what keeps the field while men sleep; and what end is awaited. Under
% each, the authorities of the different interpretations are set side by side
% at the text they share. A short comparison of the four senses closes it.
% Every claim, attribution and qualification is the expansive study's; the
% connecting prose is written for this companion.

\runninghead{Detailed Commentary}
\section{The Propers: Detailed Commentary}
\zlabel{triptych:concise:commentary:start}

\subsection{Why the servants may not weed}

The servants' offer is what any farmer's men would make, and the New
Testament seems to support it. St Jerome sets the parable's \latin{Sinite
utraque crescere}, ``Suffer both to grow,'' against St Paul's \latin{Auferte
malum de medio vestrum}, ``Take away the evil from among you'' (1~Cor 5:13):
if the Church is told both to remove the wicked and to leave them, she needs
to know which command holds when (\work{In Matthaeum} II, PL~26 cols.~93--94).
St Thomas Aquinas praises the servants for a zeal to destroy evil that the
Apostle commands and for waiting on their lord's order, \latin{iussu domini}
(\work{Super Mt.}\ c.~13); St John Chrysostom says their haste ``shows their
anxiety for the crop'' (\work{Hom.\ in Mt.}\ 46.1). The householder's refusal
gives its own reason, \latin{ne forte \dots\ eradicetis simul cum eis et
triticum}, lest the wheat be rooted up with the cockle, and the Fathers on
whom the first interpretation chiefly rests explain it.

Jerome reads it from the plant. Because darnel and wheat look alike in the
blade, the command teaches \latin{ne ubi quid ambiguum est, cito sententiam
proferamus; sed Deo judicii terminum reservemus}, ``that where anything is
doubtful we should not quickly pass sentence, but reserve to God the end of
judgment,'' so that what God removes at the end is ``not a suspicion of
crime, but manifest guilt.'' From \latin{ne forte} he draws a second point:
\latin{datur locus poenitentiae}, room is given for repentance, because one
corrupted today by a harmful teaching may tomorrow come to his senses.

Augustine reads it from the servants. His sermon on this Gospel speaks to
three groups in turn. To the zealous servants he speaks in the Lord's voice:
``ye see evil Christians among the good; and ye wish to root up the evil ones;
be quiet, it is not the time of harvest. That time will come, may it only find
you wheat!'' To the tares, in his own voice: ``amend yourselves before the
harvest come \dots\ they who to-day are tares, may to-morrow be wheat.'' To
the wheat, who ``sigh and groan as being few among many,'' he promises the
angels ``who can make the separation, and who cannot make mistakes,'' and he
places himself and his hearers among the servants, who were refused ``for ye
are such as may be deceived'': ``We are but men, the reapers are the
angels.'' Nor are the tares all below him: ``even in these high seats there is
both wheat, and tares,'' the high seats being the bench in the apse where the
bishop sat with his presbyters (\work{Sermo}~73.1--4). Elsewhere he calls the
waiting ``a certain distinction without separation'' (\work{Enarr.\ in Ps.}\
42.2--3).

Chrysostom reads it from what weeding would do. He takes the cockle for
heretics and asks why the servants are in the story at all: ``That He may
pronounce it wrong to slay them.'' In \latin{ne forte} he finds two reasons:
``If ye are to take up arms, and to kill the heretics, many of the saints also
must needs be overthrown with them; or that of the very tares it is likely
that many may change and become wheat.'' So the three agree that the cockle
may change and that the separation belongs to the harvest. They name the
servants differently. Jerome says they are angels, \latin{ne alios accipias
quam angelos}, ``understand none other than angels'' (PL~26 col.~93);
Augustine likens himself and his hearers to them; Aquinas, noting that the
reapers are angels and the servants are not, calls the servants
\latin{homines boni}, good men; and Honorius of Autun makes them prelates who
wish to separate the bad from the good (\work{Speculum Ecclesiae}). Jerome's
own warning to the hearer therefore comes from \latin{ne forte}, not from the
servants.

\subsection{How far correction may go}

The Fathers of the first interpretation do not agree how much correction the
command allows before the harvest, and this is the one real disagreement
among the authorities of the three interpretations. Jerome reserves the
doubtful case to God and does not exclude putting away the manifest one.
Chrysostom forbids killing and allows the rest: the Lord ``doth not therefore
forbid our checking heretics, and stopping their mouths, and taking away
their freedom of speech, and breaking up their assemblies and confederacies,
but our killing and slaying them.''

Augustine gives his own position in a letter to Vincentius, a Donatist
bishop. ``Originally my opinion was,'' he writes, ``that no one should be
coerced into the unity of Christ.'' He gave it up because of examples, his own
town first, brought into Catholic unity ``by fear of the imperial edicts''
(\work{Ep.}~93.17). The penalties he defends are ``such as exile and fines,''
a ``moderate severity, or rather clemency,'' meant to admonish (93.10). In the
same letter he holds to the parable: the Church looks for the visible
separation ``only on the sea-shore, i.e.\ at the end of the world,--meanwhile
correcting as many as she can, and bearing with those whom she cannot
correct'' (93.34). He argues for the penalties from their effects, not from
the parable, and his sermon on the parable names no penalty.

Aquinas reports both positions. In the \work{Summa} he notes that some have
read the command as forbidding ``not the excommunication but the slaying of
heretics, as appears from the words of Chrysostom,'' and that Augustine
changed his mind. He then says that the meaning of ``Suffer both to grow until
the harvest'' ``must be gathered from those which precede,'' the words
\latin{ne forte}, and takes from Augustine's work against Parmenian the rule
that where uprooting is not to be feared, as when a man's crime is so public
and so hated that he has no defenders, ``the severity of discipline should
not slacken'' (II-II q.~10 a.~8 ad~1). In the next question he concludes that
heretics deserve death as well as excommunication, and his grounds are not
the parable. On their side he gives the sin, ``for it is a much graver matter
to corrupt the faith which quickens the soul, than to forge money, which
supports temporal life.'' On the Church's side he gives mercy, which condemns
``not at once, but `after the first and second admonition,''' and then, ``no
longer hoping for his conversion,'' excommunicates the stubborn heretic ``and
furthermore delivers him to the secular tribunal to be exterminated thereby
from the world by death'' (q.~11 a.~3). He meets the parable in answering an
objection, with the reading he has already given: the command ``is to be
understood as referring to the case when the cockle cannot be plucked up
without plucking up the wheat'' (ad~3). The letter to Vincentius that he cites
accepts exile and fines and does not ask for death. All four hold the ground
the parable states, the safety of the wheat and the possible conversion of the
cockle; they differ over what that ground forbids before the harvest.

\subsection{What the servants do while both grow}

Paul's list in the Epistle ends in \latin{patientiam}, and his next words say
how patience acts: \latin{supportantes invicem, et donantes vobismetipsis},
bearing with one another and forgiving one another. Chrysostom explains the
first as ``passing things over,'' and reads ``even as Christ forgave you'' as
an obligation to love the offender too (\work{Hom.\ in Col.}\ 8). Aquinas says
why bearing with others is needed at all: a man may keep to justice himself
and still find that \latin{aliorum mores sunt ei importabiles}, the ways of
others are unbearable to him, as Lot was tormented living among the Sodomites
(2~Pet 2:8). To forgive, he says, is to hold no rancour and procure no harm,
\latin{sed quando necessitas puniendi est, tunc puniendus est}, ``but when
there is need to punish, then he is to be punished'' (\work{Super Col.}\ c.~3
lect.~3).

St Rabanus Maurus, on the parable, turns from the Church's judgment of others
to each man's answer to injury. The householder knew an enemy had done it,
and did not avenge himself but waited; so we, injured by our enemies,
\latin{non statim ad vindictam prorumpere, sed magis ad virtutem patientiae
recurrere debemus}, ``ought not to break out at once into revenge, but rather
to have recourse to the virtue of patience,'' and keep that long-suffering
\latin{usque ad finem vitae}, to the end of life (\work{In Matthaeum} IV,
PL~107 cols.~947--948). Neither Rabanus nor Aquinas mentions the other's
text, but both describe one act, bearing with others instead of repaying or
removing them, and Paul names that act in \latin{supportantes invicem}.

The Introit comes from the letter in which Jeremias tells the exiles to pray
for the city that holds them captive (Jer 29:7). Jerome reads that command of
the Church in this world: \latin{orare praecepit pro inimicis nostris \dots\ ut non simus
nostra tantum salute contenti, sed inimicorum quoque quaeramus salutem}, ``he
has commanded us to pray for our enemies \dots\ so that we may not be content
with our own salvation alone, but seek our enemies' salvation as well''
(\work{In Hieremiam} V, PL~24 cols.~891--892). Servants who may not weed the
field can still make that prayer for the cockle, and ask for the change
Augustine asks of the tares. St Bede, reading the Communion's promise with
the verse that follows it in Mark, sets a condition for those not yet
perfect: they will obtain what they ask, ``provided that they first forgive
others who sin against them'' (\work{In Marci evangelium expositio} III, on
11:22--25). That condition is the forgiveness Paul asks for in
\latin{donantes vobismetipsis}. And since Augustine's answer to the question
who cries from the depths is ``a sinner,'' those who sing the cry and would
weed the field are themselves among those who need mercy; in the Secret they
ask God to absolve \latin{delicta nostra}, their own sins.

\subsection{Why a God who has promised must be asked}

The second interpretation starts from the Introit, where God promises and
commands at once, and the exiles are told to call on him after the return
has been promised and before it comes. Of the parable it uses only the term
the householder sets: both are to grow until the harvest, when he will tell
the reapers what to gather (Mt 13:30). Jerome's commentary begins from the
false prophets, who claimed to know what God would do; Jeremias answers with
God's ``I know the thoughts that I think,'' and Jerome concludes,
\latin{Futurorum igitur scientia soli Deo competit}, ``knowledge of things to
come therefore belongs to God alone.'' The thoughts of peace are God's own,
known to him before anyone prays. On ``you shall call upon me \dots\ and I will hear you''
he writes that the Lord could have fulfilled his promise \latin{absque
invocatione et oratione captivorum}, without the captives' prayer, \latin{sed
hortatur eos ad preces, ut quod promissum est, mereantur accipere}, ``but he
urges them to pray, so that they may be worthy to receive what has been
promised''; and on ``you shall seek me, and shall find me'' (v.~13) he sets
the Lord's ``Ask, and you shall receive'' (Mt 7:7; \work{In Hieremiam} V,
PL~24 col.~893). The words the antiphon leaves out, \latin{finem et
patientiam}, he explains as the end of the captivity, and patience in present
labours or hope of what is to come. The Epistle asks the Colossians to put on
\latin{patientiam}, and earlier in the chapter Paul has told them that their
life ``is hid with Christ in God'' until Christ appears (Col 3:3--4). They too
hold a promise not yet fulfilled, so Jerome's gloss describes their patience
as well as the exiles'.

St Hilary of Poitiers and Chrysostom, on the psalm the Mass sings at the
Alleluia and the Offertory, speak of prayer from opposite sides. Hilary asks whether ``hear my voice'' and ``let thy ears be attentive''
say the same thing, and answers from God's side: God has no body, his ears
are his power to hear, and he is everywhere present, hearing, seeing and
acting. Yet he is to be prayed to: \latin{Naturae suae est ut audiat; sed
fidei nostrae est ut precemur audiri}, ``It is of his nature that he hears;
but it is of our faith that we pray to be heard'' (\work{Tractatus in Ps.}\
129.3, PL~9 cols.~719--720). Chrysostom answers from the side of the one who
prays. The psalmist says first that he cried and only then asks to be heard,
and it is not possible, Chrysostom says, simply to obtain from God unless
one's own part has been supplied (\work{Expositio in Ps.}\ 129, PG~55).
Neither speaks of a captivity or a promise. Their answers bear on Jerome's
question because the same Mass sings God's command, \latin{invocabitis me},
and the people's cry, \latin{De profundis}: those who cry from the depths are
the people to whom the command was given, and their cry is the call it asks
of them.

All three interpretations use the Communion. In the second, its promise
answers the Introit's: what God promised the exiles who call on him, Christ
promises to whoever asks in faith. Jerome draws a comparable link himself, between the verse the antiphon
omits and ``Ask, and you shall receive''; he does not mention the saying in
Mark. In the first, Bede's reading of the next verse makes forgiveness the
condition of being heard. In the third, Bede's reading of the verse before
names the mountain: it is the devil, called a mountain for the pride by which
he lifts himself against God, and it is cast into the sea \latin{cum
praedicantibus verbum doctoribus sanctis immundus Spiritus ab eorum corde
repellitur}, ``when, as holy teachers preach the word, the unclean spirit is
driven out of the hearts'' of those ordained to life; the perfect in faith
move such mountains \latin{orando, vel etiam jubendo}, by praying or even by
commanding (on 11:23--25). Since the Lord names the devil as the enemy who
sowed the cockle (Mt 13:39), the unclean spirit that Bede says is driven from
the heart when the word is preached, and that the perfect in faith move by
praying, is the parable's enemy, and the Communion makes its promise to that
faith.

\subsection{Thanks in affliction, and a family that leans on grace}

The Gradual gives thanks in a psalm that goes on to complain of being cast
off. Augustine, who takes its past tense for a certain future, finds his
proof in the Gradual's own second verse, ``In God we will glory'': ``Observe
how he intermingles words expressive of a future time''; only then does he ask
``what is our state at present?'' and answer it with the lament
(\work{Enarr.\ in Ps.}\ 43.7). Cassiodorus agrees, \latin{Confidentia fidei
futura pro praeteritis dicunt}, ``by the confidence of faith they speak of
things to come as past'' (\work{Expositio Psalmorum}, on Ps 43:8). St Ambrose
puts the deliverance both now and later: the enemies are confounded here,
\latin{sed multo amplius confundentur; cum viderint sanctos Domini
resurgentes}, ``but they will be confounded far more when they see the Lord's
saints rising'' (\work{Enarr.\ in Ps.}\ 43.21--22). St Robert Bellarmine reads
the verses as a memory that grounds trust: ``as often as I have been rescued
from any danger, you `have saved us from those that afflict us''' (\work{Explanatio
in Ps.}\ 43, on vv.~8--9). Chrysostom, at the verse before, tells his hearer,
when he fights the devil, to say that he trusts not in his arms or strength
but in God's mercy, and takes \latin{tota die}, all the day, for the whole of
life, so that we glory in God when he forsakes as well as when he helps
(\work{Expositio in Ps.}\ 43, §6). On every one of these answers the thanks are
sung by people still afflicted, before the deliverance is complete.

The interpretations read the Collect beside this psalm in three ways. In the
second, the family that \latin{in sola spe gratiae caelestis innititur} says
of itself what Bellarmine's psalmist and Chrysostom's hearer say, that its
trust rests on God and on nothing of its own, and its petition for protection
is the prayer of the people whose thanks the Gradual sings. In the first,
Aquinas, on the same psalm, says the saints are afflicted not only in body but
\latin{spiritualiter, quando peccant illi qui habitant apud eos}, ``spiritually,
when those who live among them sin,'' and \latin{ab his afflictionibus
liberantur a Deo}, ``from these afflictions they are freed by God''
(\work{Super Ps.}\ 43 n.~3). That is the wheat's state beside the cockle; since
the servants may not weed, the anxiety for the crop that Chrysostom saw in
their haste is left to the Collect's petition to God, who in Aquinas's
sentence is the one who frees. In the third, Chrysostom's counsel against the
devil, whom the Lord names as the enemy who sowed the cockle (Mt 13:39), is to
trust God's mercy and not one's own strength, and natural sleep, he says,
cannot be avoided; those who keep the field therefore need a keeper
themselves, and the Collect asks God to keep his family \latin{continua
pietate}, without a break.

\subsection{While men slept: the word, the peace and the song in the heart}

The first and third interpretations use the same verses of the parable and
different words of them: the first the servants' question and the
householder's answer, the third \latin{cum dormirent homines}, ``while men
were asleep,'' which the answer never mentions. Jerome, who makes the sleepers
the Church's teachers, draws a rule for them: \latin{non dormiat, qui
Ecclesiae praepositus est}, ``let him who is set over the Church not sleep''
(PL~26 col.~93). Chrysostom sets this sleep beside the wayside, the rock and
the thorns of the sower, ``so that there is need of continual watchfulness,''
and meets the obvious objection: ``as to natural sleep, it is not possible;
but as to that of our moral faculty, it is possible'' (\work{Hom.\ in Mt.}\
46.1). Rabanus puts the field in each person, who must beware the enemy's
seed, ``lest while we grow torpid in the sleep of sloth we lose the fruit of
our labour'' (PL~107 col.~947). Aquinas holds to the Lord's gloss: in this
parable, unlike the sower's, the seed is the man himself, \latin{unde non est
alia facienda expositio ab ea quam Dominus fecit}, ``so no other explanation
is to be given than the one the Lord gave'' (\work{Super Mt.}\ c.~13). The
word Paul asks to dwell in the heart is therefore not the seed of this
parable.

Chrysostom's homily on Paul's verse gives it another use. He presses both
words of \latin{Verbum Christi habitet in vobis abundanter}, ``dwell'' and
``richly,'' tells those of his hearers who have wife and children to keep at
least the New Testament at home like a chest of medicines, and names the cause
of the trouble: ``This is the cause of all evils, the not knowing the
Scriptures. We go into battle without arms, and how ought we to come off
safe?'' (\work{Hom.\ in Col.}\ 9). He does not connect this homily with the
one on the sleepers. In both, though, he describes an enemy at work against
people with no defence ready, and in this one he names the defence; so the
sleep he warns of at the parable has its remedy in the word he urges here.
Aquinas says what the word gives: wisdom, with three uses,
\latin{instructionis, devotionis et directionis}, to instruct, to move to
devotion and to direct action. Some have none of the word, and Paul says
\latin{Habitet}; others are content with a little, \latin{sed Apostolus vult
quod habeamus multum}, ``but the Apostle wants us to have much,'' and he says
\latin{abundanter} (\work{Super Col.}\ c.~3 lect.~3). The third use Aquinas
finds in verse~17, everything done in the name of the Lord Jesus; the
Secret's second petition asks God to give that direction himself, to hearts
that sway, \latin{nutantia corda tu dirigas}.

At verse~15 the Greek and the Latin give two pictures of the peace.
Chrysostom, whose text says ``let it act as umpire,'' sets a contest in the
heart: when a man is insulted unjustly ``two thoughts, the one bidding him to
revenge, the other to endure,'' wrestle, and if the peace of God stands as
umpire ``it bestows the prize on that which bids endure \dots\ by persuading
him that God is Peace, that He hath made peace with us'' (\work{Hom.\ in
Col.}\ 8). Theodoret and Theophylact, on the same Greek, keep the umpire and
the prize. Aquinas, reading \latin{exsultet}, says that joy follows from peace
and that the peace is the one Christ made between God and man. Chrysostom's
contest is the choice Rabanus sets before the injured man, revenge or
patience; the first interpretation asks what the injured man owes the
offender, the third what settles the contest within him. The Introit gives
God's own words, \latin{Ego cogito cogitationes pacis, et non afflictionis},
``I think thoughts of peace, and not of affliction''; a man who believes them believes what Chrysostom's umpire would persuade him
of, that God is for peace with him.

Paul's ``singing in grace in your hearts to God'' Chrysostom reads as ``Not
simply with the mouth \dots\ but with heedfulness'' (\work{Hom.\ in Col.}\
9). Aquinas says the canticles are sung \latin{in cordibus, et non in labiis
tantum}, in the heart and not on the lips only, and says what the Church's
singing is for: \latin{Sunt autem cantica ecclesiae cordis principaliter, sed
oris sunt ut excitetur canticum cordis}, ``The songs of the Church are
principally of the heart; but they are of the mouth so that the song of the
heart may be roused.'' He says it of the Church's songs in general, and so it
holds of the psalms this Mass sings. At Psalm~129 the second interpretation
took Hilary's sentence on faith; the third takes his picture of the psalmist
praying \latin{non labiis, sed secreto \dots\ cordis}, not with his lips but
in the secret of the heart (\work{Tractatus in Ps.}\ 129.2). Chrysostom says
that \latin{de profundis} means from the deepest heart, where prayer stands
rooted like a tree against the devil's assaults, and Bellarmine speaks of
``two depths, the depth of misery and the depth of his heart'': one who does
not reflect on the depth in which he lies ``has no wish to rise out of it''
(\work{Explanatio in Ps.}\ 129, vv.~1--2). That man, like Augustine's sinners
who are too deep even to cry, fits Chrysostom's account of the sleeper,
asleep in his moral faculty.

\subsection{The harvest, the return and the pledge}

Each interpretation sets the present before a fulfilment not yet given: the
harvest in the first and third, the return of the exiles in the second. In
the first, the harvest ends the forbearance. Rabanus's householder, \latin{qui
ante per longanimitatem patienter sustinuit, modo per justitiam meritis
condigna rependit}, ``who before bore patiently through long-suffering, now
repays by justice what the deserts merit''; and Chrysostom, who would spare
heretics so long as they stand by the wheat, says that when they have
departed ``having profited nothing, then of necessity the inexorable
punishment will overtake them.'' Augustine's wish for the servants, ``may it
only find you wheat!'', leaves even them unsure how the harvest will find
them, and the parable's end for the wheat is the barn; so the Postcommunion's
prayer for the effect of salvation is, in the first interpretation, a prayer
to be found wheat at the harvest. In the third the end is the end of the
watch: Chrysostom quotes of the sleepers ``He that endureth to the end, the
same shall be saved,'' and the Postcommunion asks for the effect of that
salvation, whose pledge has been received.

In the second the end is the return. Jerome reads the letter anagogically:
``we remain in the confusion of this age until we are found worthy to receive
the rest of the number seven,'' and then God fulfils his promise, \latin{et
reducamur ad locum nostrum Ecclesiam}, ``and we are brought back to our own
place, the Church.'' Augustine explains the past tense of the Introit's verse
from the psalm's next lines, where the psalmist, having said that God has
turned away his anger, prays ``Convert us'': ``These things I speak of as
done, because I see them about to be done: but because they are not yet done,
I pray that they may come'' (\work{Enarr.\ in Ps.}\ 84.4). The Postcommunion
prays in the same way, for what is pledged and not yet done. Bellarmine
reads the verse of a deliverance that ``will be thoroughly accomplished in the
resurrection only'' (\work{Explanatio in Ps.}\ 84, on v.~2). Three medieval
commentators on the Mass read its chants in order. Rupert of Deutz hears in
them thanks for the future salvation of the remnant of Israel, and Sicard of
Cremona the Lord's announcement that the remnant will be saved. For Rupert and
William Durandus the Offertory is the plea of Joseph's brothers, ``we
beseech thee to forget the wickedness of thy brethren'' (Gen 50:17), and the
Communion the answer of the true Joseph, Christ, \latin{Nolite timere}, ``Fear
not'' (Gen 50:19; Rupert, \work{De divinis officiis} XII.23); Durandus gives the
reason the Communion follows, \latin{quia liberalis est Dominus, et cito
dimittit}, ``because the Lord is generous and quickly forgives''
(\work{Rationale} VI.140); and Sicard has the Communion assure the remnant
again of salvation (\work{Mitrale} VI.23).

\subsection{The four senses compared}

The literal sense of the first interpretation starts from the householder's
refusal and Paul's bearing and forgiving; of the second, from Jeremias's
letter and the psalms the people sing; of the third, from the sleepers and
Paul's indwelling word. The
allegory of the first is Christ, who will judge through angels who cannot
err, and a Church in which wheat and cockle stand together in the high seats
as among the laity (Augustine); of the second, Christ turning away the
captivity of sin, the martyrs' voice in the Gradual, and a propitiation that
is the sacrifice offered for us (Augustine); of the third, the word of Christ
dwelling richly in each member of the Church (Chrysostom, Aquinas) and the
devil driven from the hearts of those ordained to life (Bede). The moral sense
of the first turns to others: judge slowly (Jerome, Augustine), seek
the conversion of the erring (Augustine, Chrysostom), bear injury (Rabanus),
and forgive in order to be heard (Bede);
of the second, to prayer within a promise, made because God urges it
(Jerome), as faith's part (Hilary), and as a sinner's (Augustine); of the
third, to oneself: stay awake in the will, keep the Scriptures and let
peace settle the contest within (Chrysostom), and pray from the heart
(Chrysostom, Aquinas, Hilary). The anagogy of the first is the
harvest, where the householder repays by justice (Rabanus); of the second, the
return ``to our own place, the Church'' (Jerome) and a deliverance
accomplished in the resurrection (Bellarmine), when the enemies are confounded
far more at the rising of the saints (Ambrose); of the third, salvation for
those who endure to the end, and canticles referred \latin{ad promissa
aeterna}, to the eternal promises (Aquinas). None of the three needs another's
conclusion, and none contradicts another.
